\section{Introduction}
Many predictors output a \emph{composition}: a non-negative vector that sums to one, such as a softmax probability vector, a vector of topic proportions or an annotator vote distribution. A user of such a prediction needs more than a point estimate; they need to know whether the attached uncertainty statement can be trusted. Split conformal prediction offers a finite-sample marginal coverage guarantee under exchangeability for any score function \citep{angelopoulos2021gentle}, and has been applied to label sets of classifiers \citep{angelopoulos2021raps,ding2023clustered}, to batch group-conditional guarantees \citep{jung2023batch} and, for composition-valued responses, to regions of the simplex \citep{javanmardi2024credal}. What is missing is an agreed, reproducible way to \emph{evaluate} such methods when the response lives on the simplex: the efficiency measure of a region of the simplex, the strata used to probe conditional coverage, the number of calibration points at which these diagnostics can be trusted, and which of the observed differences between scores are real.

We therefore treat evaluation itself as the object of study and ask: \emph{for composition-valued predictions, which parts of conformal evaluation are score-free guarantees, and which depend on the geometry of the score, the calibration scheme and the sample size?} We study a response $y\in\simplex$ and a regression $x\mapsto \hat p(x)\in\simplex$, and compare split-calibrated regions built from seven scores (Euclidean, total variation, Hellinger, KL, Aitchison/CLR, Linf and a Dirichlet negative log-likelihood), in plain, locally normalised and Mondrian variants, plus CV+. A synthetic oracle with a known conditional law provides ground truth for conditional coverage; CIFAR-10 softmax outputs and CIFAR-10H vote proportions give a real label-set bridge and a real composition-valued response.

\paragraph{Contributions.} (i) A concrete evaluation protocol for composition-valued conformal regions, with explicit choices of size measure, strata and uncertainty quantification (Section~\ref{sec:protocol}). (ii) An exact check, per score and family, of the finite-sample coverage law of split conformal prediction (Section~\ref{sec:theory}, Section~\ref{sec:exp}). (iii) An empirical comparison of scores and calibration schemes, with confidence intervals over repeated calibration sets and data seeds, and an oracle measurement of conditional coverage (Section~\ref{sec:exp}). (iv) A null simulation that quantifies how large a difference in binned conditional-coverage metrics must be to be interpretable at realistic sample sizes. (v) An explicit statement of where the conclusions hold and where they must not be extrapolated (Section~\ref{sec:disc}). We claim no new coverage guarantee and no conditional guarantee; our novelty claim is limited to the literature we read (Section~\ref{sec:related}).
